\section{From crossing decompositions to nonpositive curvature}
\label{sec:crossing}

The geodesic extension of Section~\ref{sec:geodesic} gives convexity of
wire distances. To obtain nonpositive curvature, we need a quadratic
bound on those distances. Its algebraic source is a decomposition of two
wires into two tensors that factor across different cuts of the four
endpoints. The cost of the internal wire on one cut can be traded
against its cost on the other. We use this tradeoff first to prove the
parallelogram law for normed planes realized by wire distances and then
to prove the squared midpoint inequality in every geodesic spectral
space.

\subsection{The integral crossing identity}\label{app:crossing}

The crossing construction originates in work of Anthropic~\cite{Anthropic}.
Khoruzhii~\cite{Khoruzhii} gives an integral construction in which two
products of fermionic operators factor through states below the same
energy bound. Extracting a shifted coefficient changes the two bounds
by equal and opposite amounts.
\begin{theorem}[Shifted crossing]\label{thm:crossing}
Let $t,u>0$, put $S=t^2+u^2$, and take $-S<\delta<S$. There are
integers $d_X,d_Y$, depending on $n$, with
$\log_2 d_X=nt+o(n)$ and $\log_2 d_Y=nu+o(n)$, for which the
two-wire tensor has the ordinary sum decomposition
\[
 E_{X_+X_-}(d_X)\,E_{Y_+Y_-}(d_Y)=T_++T_-.
\]
On the cuts $X_+Y_+\mid X_-Y_-$ and $X_+Y_-\mid X_-Y_+$ respectively,
$T_+$ and $T_-$ have flattening ranks at most
\[
  2^{n\sqrt{S+\delta}+o(n)}
  \quad\text{and}\quad
  2^{n\sqrt{S-\delta}+o(n)} .
\]
One can choose the tensors and their two matrix factorizations over
$\mathbb Z$, so the same choices work over every field.
\end{theorem}

The two terminal pairs $(X_+,X_-)$ and $(Y_+,Y_-)$ carry the original
wires. Each summand can be represented by joining two local factors
with a single internal wire, as in Figure~\ref{fig:crossing-cuts}.
In the integral construction, the internal wire records the intermediate
state of a product of operators together with auxiliary summation indices.
There are polynomially many choices of the auxiliary indices at the given
energy cutoff. Reversing the order of the operators exchanges the two cuts. The shift $\delta$ lets us choose how much dimension to
allow on each internal wire.

\begin{figure}[htbp]
\centering
\begin{tikzpicture}[x=1cm,y=1cm,line cap=round,
    terminal/.style={circle,fill=ink,inner sep=1.6pt},
    factor/.style={rectangle,draw=ink,fill=white,inner sep=2.5pt},
    wire/.style={line width=.7pt,draw=ink},
    internal/.style={line width=1.2pt,draw=sea},font=\small]
\foreach \a in {0,6.1}{
  \begin{scope}[xshift=\a cm]
    \node[terminal,label=above:$X_+$] (xp) at (-1.55,.68) {};
    \node[terminal,label=above:$Y_-$] (ym) at (1.55,.68) {};
    \node[terminal,label=below:$Y_+$] (yp) at (-1.55,-.68) {};
    \node[terminal,label=below:$X_-$] (xm) at (1.55,-.68) {};
    \ifdim\a pt=0pt
      \node[factor] (a) at (-.60,0) {};
      \node[factor] (b) at (.60,0) {};
      \draw[wire] (xp)--(a)--(yp) (ym)--(b)--(xm);
      \draw[internal] (a)--node[above] {$r_+$} (b);
      \node at (0,-1.45) {$T_+$};
    \else
      \node[factor] (a) at (0,.51) {};
      \node[factor] (b) at (0,-.51) {};
      \draw[wire] (xp)--(a)--(ym) (yp)--(b)--(xm);
      \draw[internal] (a)--node[right] {$r_-$} (b);
      \node at (0,-1.45) {$T_-$};
    \fi
  \end{scope}
}
\end{tikzpicture}
\caption{The two factorizations of Theorem~\ref{thm:crossing}.
In each summand the two squares are local tensors. The wire joining
them has dimension equal to the flattening rank $r_+$ or $r_-$ across
the corresponding cut. The dimensions on the four outer wires agree in
both summands.}\label{fig:crossing-cuts}
\end{figure}

\begin{proof}[Proof of Theorem~\ref{thm:crossing}]
The construction has three steps. We identify an integral pairing
of operators, factor it through states with two different energy
bounds, and count the states allowed by each bound.

\smallskip\noindent\emph{The pairing.}
Use the exterior algebra over $\mathbb Z$ on generators indexed by
$i\ge 0$. Its basis states are finite subsets of these indices,
written in increasing order. Let $c_i$ wedge the $i$th generator on
the left and let $a_i$ be its signed deletion. Creation operators
anticommute with one another, as do deletion operators, and
$a_i c_j+c_j a_i=\mathbf 1_{i=j}\mathrm{Id}$. Give a basis state $I$
energy $h(I)=\sum_{i\in I}(i+1/2)$. For background on the free fermion and vertex superalgebras see \cite{KacWang}. Introduce the fields
\[
 \psi(z)=\sum_{i\ge0}(c_i z^i+a_i z^{-i-1}),
 \qquad \psi_j(z)=\partial^{(j)}_z\psi(z),
 \qquad Y_I(z)={:}\prod_{i\in I}\psi_i(z){:}.
\]
Here $\partial^{(j)}_z z^b=\binom bj z^{b-j}$ for every integer $b$,
using the coefficient of $v^j$ in $(z+v)^b$. The product is ordered
by increasing $i$. The colons instruct us to expand and move creations
to the left of deletions using permutation signs only. Such a normal
product is skew under interchange of two fields. A coefficient $[z^b]Y_I$
shifts energy by $b+h(I)$. Indeed each summand of a differentiated
field shifts energy by its exponent plus the index of differentiation
plus $1/2$.

For equally large nonempty sets $I,J$, put $m=|I|=|J|$, $h=h(I)$,
$h'=h(J)$ and $L=h+h'-1$, a nonnegative integer. We claim
\begin{multline}\label{eq:crossing-field}
 (z-w)^L\big(Y_I(z)Y_J(w)-(-1)^{m^2}Y_J(w)Y_I(z)\big)
   \\= C_{IJ}\sum_{s\in\mathbb Z}z^{-s-1}w^s\,\mathrm{Id},
\end{multline}
where
\[
 C_{IJ}=(-1)^{\binom m2+\sum_{i\in I}i}
       \det\left[\binom{i+j}{i}\right]_{i\in I,\,j\in J}.
\]
To see this, normal order the product of the two blocks. When a
deletion from the first block meets a creation from the second, the
anticommutation relation permits either a signed interchange or a
scalar contraction. Summing the scalar for fields $i,j$ gives the
expansion in $w/z$ of
\[
  \partial^{(i)}_z\partial^{(j)}_w\frac{1}{z-w}
   =\frac{(-1)^i\binom{i+j}{i}}{(z-w)^{i+j+1}},
\]
since the undifferentiated contraction is
$\sum_{r\ge0}z^{-r-1}w^r$. Thus the expansion is indexed by partial
matchings between the blocks. Each summand is the product of the
matched scalars times the normal product of the unmatched fields.
Its sign is that of the permutation listing the matched pairs first,
in first-field/second-field order within each pair, followed by the
unmatched fields in their original order. In fact the sign for a
contraction is the sign of bringing the two operators together.
Moving the resulting pair as a block incurs no further sign.

For the reversed product, use the same ordering of pairs and the
same final order of the unmatched fields. Interchanging the original
blocks contributes $(-1)^{m^2}$. Reversing the orientation of any
pair contributes a further minus sign in the permutation, precisely
canceled by reversing its rational kernel, since
$(-1)^j/(w-z)^{i+j+1}=(-1)^{i+1}/(z-w)^{i+j+1}$. Consequently after
multiplication by $(-1)^{m^2}$ the reversed summands expand the same
rational kernels in $z/w$. For an incomplete matching an unmatched
field in each block accounts for at least one unit of missing pole
order. Its pole order is therefore at most $L$, and multiplication
by $(z-w)^L$ makes the two expansions equal. For a full matching the
pole order is $L+1$. Interleaving the two increasing lists contributes
$(-1)^{\binom m2}$, permuting partners contributes the determinant
sign, and the kernels contribute the remaining row signs in $C_{IJ}$.
The two expansions of $1/(z-w)$ differ by
$\sum_{s\in\mathbb Z}z^{-s-1}w^s$. This proves
\eqref{eq:crossing-field}.

We justify the formal products coefficientwise. Truncate the modes at
$N$ and first
perform the two normal-ordering expansions with their truncated scalar
kernels. Fix input and output states and a monomial $z^b w^{b'}$.
In the first product the intermediate energy is the input energy plus
$b'+h'$. In the reversed product the increment is $b+h$. There are
only finitely many states at either energy. Each individual normal
block between specified states uses deletion modes from its input
and creation modes from its output. Hence its matrix coefficient
stabilizes as a Laurent polynomial. The same argument applies to
the unmatched normal product, even in two variables. For each of its
finitely many monomials, prescribing the total $w$-power bounds all
contracted modes in the first expansion: a creation in mode $r$
differentiated $j$ times has $w$-power $r-j\ge0$ whenever its
coefficient is nonzero. The reversed expansion uses the $z$-power in
the same way. Both Wick expansions therefore pass coefficientwise to
the infinite-mode limit, where the rational-kernel cancellation above
applies. Multiplication by the fixed polynomial $(z-w)^L$ preserves
these limits. All coefficients are integers, including the divided
derivatives, so reduction modulo $2$ also preserves the identity.

\smallskip\noindent\emph{The factorizations.}
Let $D(H)$ count the states of energy at most $H$, with $D(H)=0$ for
$H<0$. Fix $H,E\ge1$ and an integer $s$. Use nonempty $I,J$ of
energy at most $H$, and input and output states $\beta,\alpha$ of
energy at most $E$. For a block of given $m,h,h'$ take
$[z^{-s-1}w^s]$ in \eqref{eq:crossing-field}. For the multiplier
monomial $z^{L-a}w^a$, $0\le a\le L$, the first product acts on
$\beta$ first by $[w^{s-a}]Y_J$, reaching energy at most
$E+h'+s-a\le E+H+s$. The reversed product acts first by
$[z^{-s-1-L+a}]Y_I$, reaching energy at most
$E-s-h'+a\le E+H-s$. Sum over intermediate states in each matrix
product. This writes $C_{IJ}\mathbf 1_{\alpha=\beta}$ as a sum of
two matrices, each factored over $\mathbb Z$. On the cut
$(I,\alpha)\mid(J,\beta)$ the first factors through at most
$(L+1)D(E+H+s)$ coordinates, indexed by $a$ and the intermediate
state. On the cut $(I,\beta)\mid(J,\alpha)$ the second factors
through at most $(L+1)D(E+H-s)$ coordinates. The integer coefficients
of the multiplier and the sign of the reversed product can be placed
on either factor.

Extend each block by zero and sum over $m,h,h'$. The number of blocks
and the largest $L+1$ are bounded by a polynomial in $H+1$, since
all energies are half-integers and $m\le 2H$. Define $C_{IJ}=0$ for
$|I|\ne |J|$. The resulting target is the full matrix
$C\otimes\mathrm{Id}_{D(E)}$ on the four indices. The two factorizations
have respective sizes at most
\[
  \operatorname{poly}(H+1)D(E+H+s),\qquad
  \operatorname{poly}(H+1)D(E+H-s).
\]
To turn $C$ into an identity, let $P_{i\ell}=\binom{i}{\ell}$
for nonnegative indices up to the largest mode under consideration.
Vandermonde's identity gives
$(PP^{\mathsf T})_{ij}=\sum_\ell\binom{i}{\ell}\binom{j}{\ell}
=\binom{i+j}{i}$. For fixed $m$, let $P^{[m]}$ be the matrix of
$m$-rowed minors restricted to sets of energy at most $H$ on both
indices. A minor of $P$ vanishes unless its increasing column list
is componentwise bounded by its increasing row list, since $P$ is
lower triangular. Thus restricting to this energy downset discards
no term of Cauchy--Binet, and the corresponding block of $C$ is
$P^{[m]}(P^{[m]})^{\mathsf T}$ up to row signs. In lexicographic
order $P^{[m]}$ is unit lower triangular. Every block is therefore
unimodular, and applying $C^{-1}$ at the $I$ index gives the
identity of dimension $D(H)-1$ without increasing either factorization
size.

\smallskip\noindent\emph{Counting the intermediate states.}
Write $\ln$ for the
natural logarithm. With
$A=\int_0^\infty\ln(1+e^{-x})\,dx>0$, the partition function is
\[
  \mathcal Z(\lambda)=\prod_{i\ge0}(1+e^{-\lambda(i+1/2)}),
  \qquad \ln\mathcal Z(\lambda)\sim A/\lambda
  \quad(\lambda\downarrow0).
\]
Under probabilities proportional to $e^{-\lambda h(I)}$ the occupation
variables are independent, with means $1/(e^{\lambda(i+1/2)}+1)$.
The mean energy is asymptotic to
$\lambda^{-2}\int_0^\infty x/(e^x+1)\,dx=A/\lambda^2$,
by integration by parts. Its variance is $O(\lambda^{-3})$ by
summing $(i+1/2)^2e^{-\lambda(i+1/2)}$. These Riemann-sum
estimates have exponentially decreasing tails. From
$D(H)\le e^{\lambda H}\mathcal Z(\lambda)$ at
$\lambda=\sqrt{A/H}$ we obtain the upper asymptotic bound. For the
lower bound take $0<\epsilon<1/2$ and
$\lambda=\sqrt{A/((1-\epsilon)H)}$. Chebyshev's inequality puts
probability $1-o(1)$ on energies between $(1-2\epsilon)H$ and $H$.
The number of states in this interval is at least
$(1-o(1))\mathcal Z(\lambda)e^{\lambda(1-2\epsilon)H}$.
Letting $\epsilon\downarrow0$ gives
\[
     \log_2 D(H)\sim c\sqrt H,\qquad c=2\sqrt A/\ln 2 .
\]
Finally take $H=(nt/c)^2$, $E=(nu/c)^2$ and
$s=\lfloor\delta n^2/c^2\rfloor$. Both shifted energy bounds are
positive for large $n$. Set $d_X=D(H)-1$, $d_Y=D(E)$ and identify
$(I,J,\alpha,\beta)$ with $(X_+,X_-,Y_+,Y_-)$. The polynomial
factors cost $o(n)$ in the logarithms, proving the theorem.
\end{proof}


\subsection{Rigidity of normed planes}

Place the four terminals in a normed plane. The two auxiliary sites in
each factorization may be placed anywhere. If the norm violates the
parallelogram law, we can place these sites so that both summands cost
strictly less at exponential scale than the original pair of wires.
Subadditivity of a spectral point rules this out.
\begin{lemma}[Norm rigidity]\label{lem:norm}
Let $V$ be a real normed plane. If each finite subset of some
nonempty open set in $V$ is realized at some scale by a spectral wire
metric over the fixed field, then the norm comes from an inner product.
\end{lemma}

\begin{proof}
Any finite subset of the plane can be put inside the open subset by a
translation and a positive dilation. Let $x,y$ be unit vectors at
which the norm is differentiable, and let $f,g$ be its derivatives
there, elements of the dual unit sphere. We prove that for $t,u>0$,
\begin{equation}\label{eq:crossing-dual}
   \|tf+ug\|_*^2+\|tf-ug\|_*^2\le 2(t^2+u^2).
\end{equation}
Suppose the inequality fails. Put $S=t^2+u^2$ and
$A_\pm=\|tf\pm ug\|_*$. Both are positive and
$A_\pm^2\le(t+u)^2\le 2S$. Thus one can choose $-S<\delta<S$
such that $S-A_-^2<\delta<A_+^2-S$. Take unit vectors $z_\pm$
on which $tf\pm ug$ attain their respective positive maxima.
Place the terminals at $x,-x,y,-y$. On the first cut one center serves
$x,y$ and the opposite center serves $-x,-y$. On the second cut the
partners change to $x,-y$ and $-x,y$. Starting both centers at the
origin gives each placement the same terminal cost. Moving them apart
shortens the terminal wires at a rate controlled by $tf+ug$ or $tf-ug$,
respectively, and requires a wire between the two centers. Choose the
centers $\epsilon z_+,-\epsilon z_+$ on the first cut and
$\epsilon z_-,-\epsilon z_-$ on the second, with $0<\epsilon<1$.

For any matrix factorization of a four-terminal tensor across a cut,
give each terminal an identity wire of its dimension to the center on
its side, and join the centers by a wire of the factorization size.
Contracting the indices at the centers against the two factors
produces the tensor by local restriction, with trivial output at
both centers. Coincident auxiliary roles may be grouped at one party.
Apply this observation to $T_+$ and $T_-$ of
Theorem~\ref{thm:crossing}. Realize the terminals and both pairs of
centers simultaneously, with wire distances $\sigma$ times their
norm distances for some $\sigma>0$. By multiplicativity, restriction
monotonicity, and the formula $F(E_{ab}(q))=q^{d(a,b)}$ of
Section~\ref{sec:spectra}, the log base two of $F(T_+)$, divided by
$2n\sigma$, is at most
\[
 t\|x-\epsilon z_+\|+u\|y-\epsilon z_+\|
                 +\epsilon\sqrt{S+\delta}+o(1).
\]
Here and below a tensor of value zero satisfies the corresponding
exponential bound directly. For $T_-$ the same expression applies
with $y$ replaced by $-y$, $z_+$ by $z_-$ and $\delta$ by $-\delta$.
As $\epsilon\downarrow0$ the leading expressions are
\[
 t+u-\epsilon\big(A_\pm-\sqrt{S\pm\delta}\big)+o(\epsilon),
\]
because the norm derivative at $-y$ is $-g$. Both linear savings
are strictly positive. Fix $\epsilon$ small enough for both, and
then fix the joint realization, independently of $n$. Subadditivity
under ordinary tensor sums now contradicts the wire formula for the
target, whose log value divided by $2n\sigma$ tends to $t+u$.
This proves \eqref{eq:crossing-dual}.

The closed convex hull of all derivatives at smooth unit vectors
is the dual unit ball. Indeed the norm is differentiable on a dense
set by finite-dimensional convexity. For any unit vector $v$, take
smooth unit vectors $v_j\to v$ with derivatives $f_j$. Then
$f_j(v)\ge f_j(v_j)-\|v-v_j\|\to1$, so these derivatives have the
same support function as the dual unit ball. Joint convexity and
continuity in $(f,g)$ now extend \eqref{eq:crossing-dual} to arbitrary
$f,g$ in that ball. Taking $t,u$ to be arbitrary dual vector norms
gives
\[
   \|a+b\|_*^2+\|a-b\|_*^2\le 2(\|a\|_*^2+\|b\|_*^2).
\]
Applying this inequality to $a+b,a-b$ proves its reverse. The
parallelogram law gives an inner product on the dual, and duality in
finite dimension gives the assertion for the original norm.
\end{proof}

\subsection{Nonpositive curvature}\label{sec:curvature}
On a spectral space, distances between corresponding points of two
geodesics are convex by \eqref{eq:wire-busemann}. Norm rigidity upgrades
this convexity to the squared midpoint bound characteristic of
$\operatorname{CAT}(0)$ spaces. The resulting extension preserves the prices of all
tensors on the original sites.
\begin{theorem}[Curvature]\label{thm:geometry}
Every complete geodesic spectral space $X$ is $\operatorname{CAT}(0)$.
In particular every spectral space embeds with prices in such a space, called below a \emph{Hadamard spectral space}.
\end{theorem}
We first examine two geodesics from a common origin whose separation
grows exactly linearly. They span a Euclidean sector, possibly
degenerate. We then rescale an arbitrary pair of geodesics at their
origin. Convexity makes the separation linear in this limit, so the
sector calculation gives an inequality for the original distances.

Fix a complete geodesic spectral space $X$.
Consider two constant-speed segments $a,b$ starting at $o$, defined on $[0,R]$, or two rays on $[0,\infty)$. Suppose $d(a(r),b(r))=cr$ throughout. On the coordinate triangle $u,v\ge0$, $u+v\le R$, or on the whole quadrant for rays, define
\[
 H(u,v)=[a(u+v),b(u+v)]_{v/(u+v)},\qquad H(0,0)=o.
\]
This map takes affine segments to geodesics with their affine parametrizations. Indeed, let $(u,v)$ be the affine interpolation of two parameter pairs, and $w$ the corresponding interpolation of their images. Comparing with the same interpolation on each boundary ray, \eqref{eq:wire-busemann} gives $d(w,a(u+v))\le cv$ and $d(w,b(u+v))\le cu$. Their sum is bounded by the distance $c(u+v)$ between these boundary points. Equality and uniqueness of geodesics force $w=H(u,v)$. The formula and \eqref{eq:wire-busemann} also give continuity, including at the origin.

The pullback distances of $H$ are given by a seminorm on the coordinate plane. To see this, let $N_p(v)$ be the speed of $H(p+hv)$ as a function of $h$ near $0$, for an interior point $p$. For any two interior points $p,q$ choose $s<1$ and an interior point $z$ with $q=(1-s)p+sz$. Comparing $q+hv$ to $q$ by interpolating from $p+hv/(1-s)$ and $p$ to the same $z$ shows $N_q(v)\le N_p(v)$. Interchanging $p,q$ gives equality. The common speed is absolutely homogeneous, and the triangle inequality on small segments gives subadditivity. The segment property and continuity extend the distance formula to the boundary. In rank two the seminorm is a norm, and $H$ embeds the interior of its parameter domain isometrically into $X$ with this norm. Lemma~\ref{lem:norm} therefore makes the norm Euclidean. In rank at most one it is already induced by a positive semidefinite quadratic form. Thus the sector, with all its interpolations, has the distances of a possibly degenerate Euclidean sector.

\begin{lemma}[Squared midpoint inequality]\label{lem:catzero}
For $x,y,z\in X$ and $m=[x,y]_{1/2}$,
\begin{equation}
 d(z,x)^2+d(z,y)^2\ge 2d(z,m)^2+\tfrac12 d(x,y)^2.
 \label{eq:cur-mid}
\end{equation}
At equality, the two sectors spanned at $m$ by the segments to $x,z$ and to $y,z$ have Euclidean distances, allowing degeneracy.
\end{lemma}

\begin{proof}
Take the three segments from $m$ to $x,y,z$, each on $[0,1]$. Take a pointed ultralimit at $m$ after rescaling distances by integers $j\to\infty$, still geodesic (use representative segments), carrying spectral structure by Section~\ref{sec:limits}. Take the limit of the three segments with domain rescaled by the same $j$, obtaining three rays (possibly constant). For two original segments $a,b$, the ratio $d(a(s),b(s))/s$ is nondecreasing by \eqref{eq:wire-busemann} and bounded by the sum of their speeds. Its limit at zero therefore gives a constant $c$ for which the equal-parameter distances of the limiting rays are $cr$.

Put $A=d(x,m)=d(y,m)$ and $B=d(z,m)$, initially both positive. The rays for $x$ and $y$ point in opposite directions on one geodesic. The preceding sector argument gives Euclidean angle cosines $u,v$ for their respective pairs with the third ray. The triangle inequality between the first two rays at parameter $T$, through the third at parameter $1$, reads
\[
 2TA\le\sqrt{(TA)^2+B^2-2TABu}+\sqrt{(TA)^2+B^2-2TABv}.
\]
Subtract $2TA$ and let $T\to\infty$ to find $u+v\le0$. Each original distance $d(z,x),d(z,y)$ is at least the corresponding limiting distance at parameter $1$, by monotonicity of the ratios. The sum of their squares is therefore at least $2(A^2+B^2)-2AB(u+v)\ge2(A^2+B^2)$, proving \eqref{eq:cur-mid}. Equality forces each original ratio at parameter $1$ to equal its limit at zero. The ratio is then constant on $(0,1]$, and the sector argument applies to each pair of original segments. If $A$ or $B$ vanishes, the assertions follow directly on a segment.
\end{proof}

Applying \eqref{eq:cur-mid} at dyadic fractions and passing to the limit gives the CAT(0) inequality in its squared-distance form
\begin{equation}
 d(z,[x,y]_s)^2\le (1-s)d(z,x)^2+s d(z,y)^2-s(1-s)d(x,y)^2 .
 \label{eq:cur-cn}
\end{equation}
The midpoint characterization \cite[Ch.~II.1]{BH} now gives the curvature
assertion of Theorem~\ref{thm:geometry}. Corollary~\ref{cor:spectral-geodesic}
gives the embedding assertion. To understand when a finite configuration
spans a Euclidean rectangle, we use the four-point consequence of
\eqref{eq:cur-mid}. For $A,B,C,D\in X$ it reads
\begin{equation}
 d(A,C)^2+d(B,D)^2
 \le d(A,B)^2+d(B,C)^2+d(C,D)^2+d(D,A)^2 .
 \label{eq:cur-quad}
\end{equation}
To derive this bound, apply \eqref{eq:cur-mid} on $AC$ with $z=B,D$, and use $d(B,D)^2\le 2d(B,m)^2+2d(m,D)^2$. At equality, $m$ is also the midpoint of $BD$, and \eqref{eq:cur-mid} is sharp on either diagonal with either remaining vertex. Consequently the two side squares incident at any vertex sum to $(d(A,C)^2+d(B,D)^2)/2$, and opposite sides have equal lengths.

\subsection{Filling rectangles and cubes}

Suppose $A,B,C,D$ have the six distances of the Euclidean rectangle with cyclic coordinates $(0,0),(a,0),(a,b),(0,b)$, where $a,b>0$. Put $E_s=[A,B]_s$, $F_s=[D,C]_s$. Quadrilateral equality gives a common diagonal midpoint $m$. Equality on $AC$ with $z=B$ gives a Euclidean sector bounded by $mA,mB$, by Lemma~\ref{lem:catzero}. Hence $d(E_s,m)$ equals its rectangle-model value. The same holds for $d(F_s,m)$. Each of $d(E_s,C)$ and $d(E_s,D)$ has the Euclidean upper bound by \eqref{eq:cur-cn} on $AB$. For the matching lower bound, apply \eqref{eq:cur-mid} with $z=E_s$ on $AC$ or $BD$, using the known distances to $A$ or $B$ and to $m$. The same argument determines $d(F_s,A)$ and $d(F_s,B)$.

By \eqref{eq:wire-busemann}, $d(E_s,F_s)\le b$. For the reverse inequality apply \eqref{eq:cur-quad} to $A,E_s,F_s,D$. Both diagonals have squared length $s^2a^2+b^2$, while the other three sides have lengths $sa,sa,b$. Finally, for $0<s<1$ the function
\[
 t\longmapsto d(E_s,F_t)^2-a^2(t-s)^2-b^2
\]
is convex by \eqref{eq:cur-cn} on subsegments of $DC$. It vanishes at $0,s,1$, hence everywhere. Continuity includes $s=0,1$. All cross distances between the two interpolated sides are therefore Euclidean.

It follows that corners in $X$ with the distances of a scaled Euclidean cube span an isometric cube. Fill the coordinates successively. At each stage, any two segments parallel to the coordinate being filled either coincide or have a Euclidean rectangle as their endpoint configuration in the previous stage. The cross-distance calculation preserves all distances at the next stage. The resulting image is compact and convex, since Euclidean segments map to the unique geodesics between their endpoints. Each specified cube point can be constructed with finitely many interpolations of the original corners.
